\documentclass[11pt]{article}
\usepackage[T1]{fontenc}
\usepackage{lmodern}
\usepackage[margin=0.8in]{geometry}
\usepackage{amsmath,amssymb,amsthm}
\usepackage{booktabs,longtable,array,enumitem}
\usepackage[hidelinks]{hyperref}
\newcommand{\Fix}{\operatorname{Fix}}
\newcommand{\Aut}{\operatorname{Aut}}
\newcommand{\Sel}{\operatorname{Sel}}
\newcommand{\Eqs}{\mathsf{EqSrc}}
\title{EQSRC-FINITE-COUNTERMODEL-ATLAS-V1\\
Finite Models and Countermodels of the Fixed-Core Selector Theorem}
\author{Candidate Constructor task overlay, RT-20260720-017}
\date{2026-07-21}

\begin{document}
\maketitle

\section{Status and exact scope}

This artifact executes only v21 work item P2-T04. It integrates four named
historical finite families under
\texttt{EQSRC-NO-NATURAL-SELECTOR-NONUNIQUENESS-THEOREM-V1} while preserving
their original identifiers, hypotheses, negative-result scope, and freeze
status. The atlas is \texttt{draft/control}, \texttt{proposal-only}, and
\texttt{source-extension data}. It neither supersedes the source artifacts nor
reopens a frozen candidate family.

The automorphisms below are structural permutations of declared finite source
data. They are not physical gauge transformations. Fixed-point and orbit
calculations are finite mathematical controls, not evidence of physical
covariance, physical admissibility, ontology adoption, general \(\Eqs\), or
Distance-to-GR progress.

\begin{verbatim}
atlas_id: EQSRC-FINITE-COUNTERMODEL-ATLAS-V1
plan_task_id: P2-T04
result: constructed_candidate
fixture_count: 12
historical_family_count: 4
direct_theorem_instance_count: 5
fixed_point_positive_control_count: 3
multiple_fixed_point_control_count: 1
guarded_non_instance_count: 3
current_ontology_derives_selector_package: false
structural_automorphisms_are_physical_gauge: false
exact_family_freezes_reopened: false
distance_to_gr_ledger_changed: false
metric_use_ledger_changed: false
physics_promotion_authorized: false
adoption_status: blocked_adoption_open_continuation
\end{verbatim}

\section{Common theorem coordinates}

Hold the P2-T01 certified-isomorphism core \(\mathcal G_\Sigma\), eligible
choice functor \(S_\chi\), structural actions, and relation map \(K\) fixed
before evaluating any desired outcome. For one represented connected
component with representative \(X\), the P2-T03 theorem gives
\[
 \Sel_\chi\;\cong\;
 \Fix_{\Aut_{\mathcal G_\Sigma}(X)}\!\bigl(S_\chi(X)\bigr).
\]
Thus an empty fixed set gives no deterministic source-natural selector, a
singleton fixed set gives one selector, and two fixed choices give two
selectors. The induced relation count is separately the cardinality of the
\(K\)-image and may be smaller than the selector count.

Each admitted atlas record fixes the following data before calculation:
\begin{enumerate}[label=(H\arabic*)]
  \item one finite source object or represented component;
  \item its eligible-choice set \(D\);
  \item a structurally declared automorphism group \(G\curvearrowright D\);
  \item the orbit decomposition and fixed set \(D^G\); and
  \item any relation map \(K:D\to\mathcal R\), kept distinct from selection.
\end{enumerate}
The three guarded noninstances fail a proposal admission rule before the
eligible selector domain is formed. They preserve fail-closed controls but do
not count as counterexamples to the P2-T03 theorem.

In the controlled vocabulary used below, a \emph{direct theorem instance}
meets the fixed-core hypotheses exactly, while a guarded noninstance does not.
Structural automorphisms are not physical gauge transformations. The atlas
also includes one explicit multiple-fixed control.

\section{Atlas summary}

\small
\begin{longtable}{p{0.21\linewidth}p{0.15\linewidth}p{0.12\linewidth}p{0.13\linewidth}p{0.29\linewidth}}
\toprule
Fixture & Status & \(|G|,|D^G|\) & Orbits & Exact conclusion \\
\midrule
\endfirsthead
\toprule
Fixture & Status & \(|G|,|D^G|\) & Orbits & Exact conclusion \\
\midrule
\endhead
\texttt{RESP-XEMPTY-SIGN} & direct & \(2,0\) & \(2\) & no invariant sign selector \\
\texttt{RESP-XEMPTY-TOKEN} & direct & \(2,0\) & \(2\) & no invariant token-map selector \\
\texttt{RESP-TAGGED} & positive & \(1,1\) & \(1\) & one supplied-tag response control \\
\texttt{ORIENTATION-GL2} & direct & \(6,0\) & \(3\) & no invariant proper line \\
\texttt{ROOTED-C2} & direct & \(2,0\) & \(2\) & no invariant singleton root \\
\texttt{ROOTED-C4} & direct & \(4,0\) & \(4\) & no root selector; two \(K\)-images \\
\texttt{ROOTED-LOADED} & positive & \(1,1\) & \(1\) & one selector after a root is supplied \\
\texttt{GRADED-CHAIN4} & positive & \(1,1\) & \(1\) & one conditional unique-minimum selector \\
\texttt{TRIVIAL-TWO} & multiple & \(1,2\) & \(1+1\) & two selectors; \(K\)-image one or two \\
\texttt{GRADED-PERIODIC} & guarded & \(4,0\) & \(4\) & noninstance: antisymmetry fails \\
\texttt{GRADED-MULTIMIN} & guarded & \(2,0\) & \(2\) & noninstance: unique minimum fails \\
\texttt{GRADED-NONGRADED} & guarded & \(1,2\) & \(1+1\) & noninstance: rank is multivalued \\
\bottomrule
\end{longtable}
\normalsize

\section{Response-tag family}

The preserved source is the Refuter stress artifact from
\texttt{RT-20260614-055}, including
\(F_{\rm toy}=\{X_+,X_{\rm flip},X_{\rm scale},X_{\rm token},X_\emptyset\}\)
and freeze label \texttt{NDCL-RESP-LC-SELECTOR-UNDERDETERMINATION}.
On \(X_\emptyset\), sign reversal acts transitively on
\(\{\varepsilon_+,\varepsilon_-\}\), and token relabeling acts transitively
on \(\{\tau_{ab},\tau_{ba}\}\). Both fixed sets are empty, so each is a
direct one-component instance of the no-selector branch.

For a tagged member such as \(X_+\), the stored
\((\varepsilon,\lambda,\tau)\) is a singleton eligible response under the
tag-preserving record. This is a fixed-point-positive control only. It confirms
conditional sufficiency of explicit tags; it does not derive those tags from
the untagged record and does not repair tag-erasure fragility.

\section{Orientation-line family}

Preserve
\texttt{CM-EQSRC-ORIENTATION-TORSOR-LINE-SELECTION-001} and
\texttt{OB-EQSRC-ORIENTATION-TORSOR-LINE-SELECTION-001} from
\texttt{RT-20260718-041}. The eligible choices are the three nonzero proper
lines of \(\mathbb F_2^2\). The six matrices of
\(\mathrm{GL}(2,2)\) act transitively on these lines, so their single orbit has
size three and the fixed set is empty. The affine group has order 24; each line
relation has eight ordered pairs, each affine stabilizer has order eight, and
their common stabilizer is the translation group of order four. These counts
are reproduced independently by the task-local validator.

The theorem therefore recovers the exact line-selection obstruction on the
declared proposal core. It does not turn structural \(\mathrm{GL}(2,2)\)
symmetry into physical gauge, nor does it prove a global source-extension
no-go.

\section{Rooted-partition family}

Preserve \texttt{CM-EQSRC-ORDERED-MOTION-ROOTED-PARTITION-001} from
\texttt{RT-20260718-045}. On the two-token directed cycle, the order-two
automorphism swaps the two eligible roots. The unique orbit has size two and
the fixed set is empty. On the directed four-cycle, rotations form \(C_4\),
act transitively on four eligible roots, and again leave no root fixed.

The four roots induce only two parity relations:
\[
 K(0)=K(2)=R_{02},\qquad K(1)=K(3)=R_{13}.
\]
The relation stabilizer has order two. Hence this historical example also
instantiates the P2-T03 warning that selector multiplicity and relation-image
multiplicity are distinct.

If a root certificate is supplied in advance, the root-preserving eligible
space is a singleton and gives one conditional selector and one relation. This
is a positive control, not a derivation of root provenance. The original
rooted-partition construction--audit--stress cycle remains locally frozen.

The validator also independently reproduces the historical four-token census:
4096 loop-free directed graphs, 1088 admitted rooted records, 557 admitted
graphs, 292 one-relation graphs, and 265 two-relation graphs.

\section{Graded-orbit family}

Preserve
\texttt{EqSrcFlowGeneratedGradedOrbitRootLaw\_src\string^cand,v1} from
\texttt{RT-20260718-047}. Its admitted four-event action chain has unique
minimum \(e_0\), ranks \(0,1,2,3\), and parity blocks
\(\{e_0,e_2\}\) and \(\{e_1,e_3\}\). The automorphism group of the ordered
chain is trivial and the eligible minimum is a singleton, so this is a
fixed-point-positive theorem control. The ordered action remains a
proposal-only source extension not derived by current ontology.

Three finite controls remain outside the admitted selector domain:
\begin{itemize}
  \item the periodic four-cycle fails antisymmetric reachability;
  \item the three-point action with minima \(a,b\) fails unique-minimum
        admission; and
  \item the five-point example with saturated path lengths two and three fails
        graded-rank admission.
\end{itemize}
Their action and choice counts are retained, but their exact status is
\emph{guarded noninstance}. The validator reproduces the original unary-action
census for carrier sizes one through four.

\section{Multiple-fixed and relation-image control}

Let a trivial group act on \(D=\{a,b\}\). Then
\(D^G=D\), so the one-component theorem gives two selectors. For an injective
map \(K_{\rm inj}\), the induced relation image has size two. For a constant
map \(K_{\rm const}\), it has size one. This is the required multiple-fixed
control and directly prevents the overread ``multiple selectors imply multiple
relations.''

\section{Provenance and excluded pseudo-symmetries}

Every fixture cites its original task and immutable object or candidate ID.
The machine fixture manifest also pins the five controlling source hashes.
No group generator depends on a file name, file order, repository path,
serialization order, registry order, validator order, agent role, handoff, or
process state. Such metadata may help retrieve a record but cannot be source
structure or selector evidence.

The atlas does not change the historical result vocabulary. Direct instances
remain scoped obstructions on their declared proposal domains. Positive
controls remain conditional on supplied proposal data. Guarded noninstances
remain fail-closed. No entry is adopted, rejected globally, or promoted.

\section{Distance-to-GR and next burden}

The active milestone remains \texttt{source\_equivalence\_eqsrc} and the
authoritative status remains \texttt{draft object exists}. The atlas
demonstrates theorem scope and preserves negative-result provenance, but it
does not derive or adopt the P2-T01 category, selector package, physical gauge
meaning, or general \(\Eqs\). It changes neither scientific ledger.

The next dependency-ready v21 item is P2-T05, the proof-assistant encoding of
the selector theorem core. P2-T06 remains dependent on both P2-T04 and P2-T05
and retains the independent Smuggling Auditor role. P2-T07 retains independent
Refuter stress. This atlas executes none of those later tasks.

\end{document}
